% !PRINTED 153
\weberpage{153}

% [restored 2026-08-06, 10x zoom strip 1/10] This line stands between the running head and the
% first display and was dropped by the transcription pass. It is not at the foot of i0177
% either, so it was lost rather than moved. Raised as an omission query by the English
% translation pass on this page, then confirmed against the scan.
Hierin ist nach \S.\,42 zu setzen
\[
\begin{aligned}
s_0 &= 3, \quad a_0 s_1 = -a_1\\
a_0^2 s_2 &= a_1^2 - 2 a_0 a_2\\
a_0^3 s_3 &= -a_1^3 + 3 a_0 a_1 a_2 - 3 a_0^2 a_3\\
a_0^4 s_4 &= a_1^4 - 4 a_0 a_1^2 a_2 + 4 a_0^2 a_1 a_3 + 2 a_0^2 a_2^2,
\end{aligned}
\]
wodurch man erhält

\weq{10}{D = a_1^2 a_2^2 + 18 a_0 a_1 a_2 a_3 - 4 a_0 a_2^3 - 4 a_1^3 a_3 - 27 a_0^2 a_3^2.}

\begin{center}
\S.\,47.

\gesperrt{Discriminanten der Formen dritter und vierter Ordnung.}
\end{center}

Bei den Functionen dritter und vierter Ordnung haben wir zur Berechnung der Discriminanten in den Auflösungen der Gleichungen dieser beiden Grade (\S.\,35, 37) ein einfaches Mittel.

Was zunächst die Gleichung dritter Ordnung betrifft, so waren die drei Wurzeln von
\weq{1}{x^3 + a x + b = 0}
in der Form dargestellt:
\weq{2}{\begin{aligned}
\alpha &= u + v\\
\beta &= \varepsilon u + \varepsilon^2 v\\
\gamma &= \varepsilon^2 u + \varepsilon v,
\end{aligned}}
worin
\weq{3}{\varepsilon = \frac{-1 + \sqrt{-3}}{2}, \quad \varepsilon^2 = \frac{-1 - \sqrt{-3}}{2}}
die complexen dritten Einheitswurzeln sind und $u$, $v$ durch \S.\,35, (8) bestimmt sind.

Aus (3) folgt
\weq{4}{\varepsilon - \varepsilon^2 = \sqrt{-3}, \quad 1 - \varepsilon = \varepsilon^2 \sqrt{-3}, \quad 1 - \varepsilon^2 = -\varepsilon \sqrt{-3}}
und aus (2)
\[
\begin{aligned}
\alpha - \beta &= \sqrt{-3}\,(\varepsilon^2 u - \varepsilon v)\\
\alpha - \gamma &= -\sqrt{-3}\,(\varepsilon u - \varepsilon^2 v)\\
\beta - \gamma &= \sqrt{-3}\,(u - v).
\end{aligned}
\]
Hieraus durch Multiplication
\weq{5}{(\alpha - \beta)(\alpha - \gamma)(\beta - \gamma) = 3 \sqrt{-3}\,(u^3 - v^3).}
